\begin{abstract}
O monitoramento da qualidade do ar em ambientes urbanos exige estimativas confiáveis de poluentes atmosféricos, como o dióxido de nitrogênio (NO\textsubscript{2}), cujos níveis estão diretamente associados a impactos à saúde pública. Sensores de baixo custo constituem uma alternativa promissora às estações de referência, porém suas medições apresentam alta variabilidade, deriva e forte multicolinearidade entre preditores, demandando calibração robusta. Este trabalho compara diferentes modelos de regressão para estimar NO\textsubscript{2} real a partir das respostas de sensores químicos do conjunto de dados \textit{Air Quality UCI}. Após um pré-processamento composto por transformação logarítmica e padronização, quatro abordagens foram avaliadas: Regressão Linear Ordinária (OLS), Ridge Regression, Partial Least Squares (PLS) e Multilayer Perceptron (MLP). A seleção de hiperparâmetros foi realizada por validação cruzada estratificada (10-fold), utilizando RMSE e R\textsuperscript{2} como métricas. Os resultados indicam que modelos supervisionados baseados em componentes (PLS) e regressão penalizada (Ridge) apresentam desempenho superior e maior estabilidade frente à multicolinearidade, alcançando elevada capacidade preditiva com menor complexidade em comparação ao MLP. Conclui-se que, para a calibração de sensores de baixo custo de NO\textsubscript{2}, modelos lineares bem regularizados fornecem excelente compromisso entre acurácia, interpretabilidade e viabilidade computacional. O código e dados utilizados estão disponíveis em \url{https://github.com/diego-nac/ica-homeworks}.
\end{abstract}


\begin{IEEEkeywords}
Air quality, Nitrogen Dioxide (NO\textsubscript{2}), low-cost sensors, regression models, Ridge regression, Partial Least Squares (PLS), Multilayer Perceptron (MLP), multicollinearity.
\end{IEEEkeywords}